\documentclass[UTF8,fontset=windows,zihao=-4]{ctexart}
\usepackage{amsmath,amssymb,graphicx,booktabs,longtable,array,geometry}
\geometry{a4paper,margin=23mm}
\setlength{\parskip}{0.35em}
\setlength{\emergencystretch}{2em}
\setlength{\tabcolsep}{3pt}
\renewcommand{\arraystretch}{1.25}
\setcounter{tocdepth}{2}
\usepackage[unicode,hidelinks]{hyperref}
\begin{document}

\begin{center}{\LARGE\bfseries 有限角光源下定日镜场的光学评价与分层参数化设计}\end{center}

\section*{摘要}

针对固定定日镜场评价、统一参数设计和逐镜参数设计三个相互衔接的问题，建立以有限镜面和光线路径为基础的热功率评价模型。根据题给太阳位置及辐照度关系确定镜面姿态，以方向能量份额衔接法向直接辐照度，联合处理入射阴影、反射遮挡与有限圆柱侧面接收。采用带能量权重的随机积分，先池化逐镜逐时原始统计，再形成条件效率和规定时点的样本平均；冻结设计后的独立批次用于数值确认。

固定1745面镜场的年样本平均热功率为 35.249 MW，单位镜面面积功率为 0.5611 kW/m\textsuperscript{2}。统一参数设计采用圆环与六扇区直边层带的有界搜索，经可行性修复得到2981面镜的R027设计：塔位为(0,\(-\)15)m，镜面尺寸6.65m\(\times\)6.65m，安装高度6m；年平均功率为 60.158 MW，单位面积功率为 0.4563379 kW/m\textsuperscript{2}。进一步固定其塔位和镜位，将逐镜参数划分为6组，得到Q3F013设计，年平均功率为 60.020 MW，单位面积功率为 0.4606881 kW/m\textsuperscript{2}。同一组新配对数据表明，后者在总面积减少 1543.202 m\textsuperscript{2}、总功率下降 139.4496 kW的同时，单位面积功率提高约 0.94995\%。其额定数值工作余量仅为 12.45 kW。

两项设计均满足所采用的60MW额定下限及汇总数值工作判据。该结论限于规定60个时点、有限搜索范围和本文物理近似；数值工作指标不包含尚未量化的光源、镜面误差与省略实体影响，不能解释为工程裕度或全局最优性保证。

\textbf{关键词：} 定日镜场；条件能量比；随机射线积分；参数化布局；独立确认

\section*{1 问题重述与总体分析}

题目在固定圆形场地、接收器及代表时点下提出三个任务：首先评价已给坐标的1745面定日镜；其次在统一尺寸和安装高度下设计塔位、镜数与镜位；最后允许每面镜的尺寸与安装高度分别取值。后两问均要求年平均总输出热功率达到60MW，并尽可能提高单位总镜面面积的年平均功率。[1]

评价与设计共享同一物理关系，但输入自由度不同。第1问只需把给定镜场映射为效率和功率，不能将后两问的额定条件施加到该固定镜场。后两问中，面积、姿态、目标距离和遮挡共同随设计改变；增大镜面或增加镜数既可能增加采光，又可能减少其他镜面的有效接收。因此，需要把完整设计展开后重新评价，不能将第1问效率按镜数或面积线性外推。

模型的关键是使几何事件、能量分母与统计总体一致。入射段未遮阴、反射段未遮挡和有效接收是相互关联的事件，不能把无条件损失率当成独立因素反复相乘。另一方面，单位面积目标是总功率与实际总面积的比值；在异面积设计中，镜数平均效率不能代替面积加权效率。以下先建立公共评价关系，再分别处理固定镜场、统一参数布局和固定布局下的分组异尺寸搜索。

\section*{2 公共假设与符号}

坐标轴依次指向东、北、上。场地圆心固定为平面原点；塔位为 \(\boldsymbol a=(a_x,a_y)\)，接收器中心为 \(\boldsymbol c_R=(a_x,a_y,80)\) m。接收器为半径3.5m、高8m的有限圆柱，只将侧面作为有效受光面，顶底面视为不透明非受光端面。镜面宽边始终水平，镜心安装高度与镜面边长中的“高度”分别记为 \(z_i\) 与 \(h_i\)。

采用365日非闰年代表日历，以3月21日为第0天，将题给当地时间直接代入附录时角关系，不另作民用时修正。每月21日的9:00、10:30、12:00、13:30、15:00形成60个规定样本。本文“年平均”均为这些时点的等权平均，不是连续全年热能、发电量或实测运行统计。

光源为轴对称、按立体角均匀的硬截止圆锥，参考半角为 \(\beta=0.00465\,\mathrm{rad}\)。该角尺度来自本地审核资料中的工程示例；示例原有过渡区与硬截止并不完全相同，本文明确采用后者作为建模近似，不将其称为所在地实测分布。[3]

镜面为零厚度、双面不透明的有限平面矩形，只计单次理想反射。反射率采用题给常数示例 \(\rho=0.92\)，不将该示例说成唯一可能值。未给出尺寸的塔身、支架、厂房不建立实体；未给出参数的坡度、跟踪、散射及镜间二次反射不另行加入。这些省略不表示相关实际损失为零。输出热功率沿题给光学接收口径，未另扣接收器吸收、再辐射、对流及换热损失。


\par\medskip\noindent\begin{minipage}{\linewidth}
\small
\textbf{表1　公共符号及单位}\par\smallskip
\begin{tabular}{>{\raggedright\arraybackslash}p{0.340000\dimexpr\linewidth-6\tabcolsep\relax}>{\raggedright\arraybackslash}p{0.490000\dimexpr\linewidth-6\tabcolsep\relax}>{\raggedright\arraybackslash}p{0.170000\dimexpr\linewidth-6\tabcolsep\relax}}
\toprule
符号 & 含义 & 单位 \\
\midrule
\(N,\boldsymbol p_i,\boldsymbol c_i\) & 实际镜数、镜心平面坐标、镜心三维坐标 & 1、m、m \\
\(w_i,h_i,z_i\) & 镜宽、镜高、镜心安装高度 & m \\
\(A_i,A_{\rm tot}\) & 单镜面积与总镜面面积 & m\textsuperscript{2} \\
\(\boldsymbol s_{0k},\boldsymbol s\) & 指向太阳中心及抽样光源的单位方向 & 1 \\
\(\boldsymbol t_i,\boldsymbol n_{ik},\boldsymbol u_{ik},\boldsymbol v_{ik}\) & 中心目标方向、法向、宽边及高边单位方向 & 1 \\
\(D_k,ST_k,\delta_k,\omega_k\) & 春分起算日数、当地时刻、赤纬、时角 & 日、h、rad、rad \\
\(\mathrm{DNI}_k,c_{ik},\tau_i\) & 参考法面辐照度、中心余弦、大气透射率 & kW/m\textsuperscript{2}、1、1 \\
\(S,B,R\) & 未遮阴、未反射遮挡、有效侧入接收指示量 & 1 \\
\(\Pi_{0,ik},\Pi_{1,ik},\Pi_{2,ik}\) & 基准反射、联合存活、有效接收三层功率 & kW \\
\(\eta_{ik},P_k,\overline q\) & 总光学效率、逐时全场热功率、年均单位面积功率 & 1、kW、kW/m\textsuperscript{2} \\
\(U_Q,\Delta q\) & 汇总量的数值工作指标、配对单位面积功率差 & 与相应量相同、kW/m\textsuperscript{2} \\
\bottomrule
\end{tabular}
\end{minipage}\par\medskip

\section*{3 共用的几何、光学与统计模型}

\subsection*{3.1 太阳方向与镜面姿态}

\par\noindent\begin{minipage}{\linewidth}
当地纬度为 \(\varphi=39.4^\circ\)，题给赤纬和时角公式可写为式（1）。太阳方向用带符号向量表示，从而保留上午和下午的东西分支。

\[
\begin{aligned}
\delta_k&=\arcsin[\sin(2\pi D_k/365)\sin(23.45^\circ)],
&\omega_k&=\frac{\pi}{12}(ST_k-12),\\
\boldsymbol s_{0k}&=\begin{pmatrix}
-\cos\delta_k\sin\omega_k\\
\sin\delta_k\cos\varphi-\cos\delta_k\cos\omega_k\sin\varphi\\
\sin\delta_k\sin\varphi+\cos\delta_k\cos\omega_k\cos\varphi
\end{pmatrix}.
\end{aligned}\tag{1}
\]
\end{minipage}\par

\par\noindent\begin{minipage}{\linewidth}
太阳高度满足 \(\sin\alpha_k=s_{0k,z}\)。海拔 \(H=3\,\mathrm{km}\) 仅用于题给辐照度公式，与安装高度及接收器高度区分：

\[
\begin{aligned}
\mathrm{DNI}_k&=1.366[a+b\exp(-c_{\rm D}/\sin\alpha_k)],\\
a&=0.4237-0.00821(6-H)^2,\\
b&=0.5055+0.00595(6.5-H)^2,\\
c_{\rm D}&=0.2711+0.01858(2.5-H)^2.
\end{aligned}\tag{2}
\]
\end{minipage}\par

\par\noindent\begin{minipage}{\linewidth}
式（2）的单位为kW/m\textsuperscript{2}，指数系数 \(c_{\rm D}\) 用下标与余弦因子区分。由镜心 \(\boldsymbol c_i=(x_i,y_i,z_i)\) 瞄准接收器中心，镜面法向取太阳中心方向与目标方向的角平分方向：

\[
\begin{aligned}
\boldsymbol t_i&=\frac{\boldsymbol c_R-\boldsymbol c_i}{\|\boldsymbol c_R-\boldsymbol c_i\|},
&\boldsymbol n_{ik}&=\frac{\boldsymbol s_{0k}+\boldsymbol t_i}{\|\boldsymbol s_{0k}+\boldsymbol t_i\|},\\
\boldsymbol u_{ik}&=\frac{(-n_{ik,y},n_{ik,x},0)}{\sqrt{n_{ik,x}^2+n_{ik,y}^2}},
&\boldsymbol v_{ik}&=\boldsymbol n_{ik}\times\boldsymbol u_{ik}.
\end{aligned}\tag{3}
\]
\end{minipage}\par

\par\noindent\begin{minipage}{\linewidth}
仅当法向两个水平分量严格为零时，规定宽边向东；近竖直姿态仍采用式（3）。有限镜面源点与理想反射方向为

\[
\begin{aligned}
\boldsymbol x&=\boldsymbol c_i+\xi\boldsymbol u_{ik}+\zeta\boldsymbol v_{ik},
&|\xi|&\le w_i/2,\quad |\zeta|\le h_i/2,\\
\boldsymbol r&=-\boldsymbol s+2(\boldsymbol s\cdot\boldsymbol n_{ik})\boldsymbol n_{ik}.
\end{aligned}\tag{4}
\]
\end{minipage}\par

其中 \(\boldsymbol s\) 指向光源，实际入射传播方向为 \(-\boldsymbol s\)。区分位置和方向可避免把太阳回查射线与入射传播方向用反。

\subsection*{3.2 有限相交与有效接收}

\par\noindent\begin{minipage}{\linewidth}
对射线 \(\boldsymbol y(t)=\boldsymbol x+t\boldsymbol d\)，矩形交点先由镜面平面求得正参数，再检查两个局部坐标的有限边界。无限平面的相交、负参数交点以及自交均不能代替有限镜面命中。接收器侧面和端面满足

\[
\begin{aligned}
\text{侧面：}\quad &(X-a_x)^2+(Y-a_y)^2=3.5^2,\quad |Z-80|\le4,\\
\text{端面：}\quad &(X-a_x)^2+(Y-a_y)^2\le3.5^2,\quad Z\in\{76,84\}.
\end{aligned}
\tag{5}
\]
\end{minipage}\par

所有相交事件按最小正参数确定首次实体接触；只有首次接触为侧面且方向向实体内部时才计入有效接收。端面先到、远侧出射、内部向外穿出和切触不视为有效外入。切触、接缝与机器容差带内的事件单列不确定状态，不扩大物理边界，不移动起点强制消除异常。接收器既可以成为入射路径障碍，又是具有专门有效面定义的接收对象，两种身份分别判断。

\subsection*{3.3 光源方向份额与投影权重}

\par\noindent\begin{minipage}{\linewidth}
令 \(\theta\) 为方向与太阳中心轴的夹角，\(\psi\) 为轴周方位角。方向能量份额相对于立体角的密度为

\[
p_\Omega(\theta,\psi)=\frac{1}{2\pi(1-\cos\beta)},
\quad 0\le\theta\le\beta,\quad0\le\psi<2\pi,
\qquad\int p_\Omega\,\mathrm d\Omega=1.
\tag{6}
\]
\end{minipage}\par

\par\noindent\begin{minipage}{\linewidth}
这里的“均匀”不是极角等概率，也不是未经转换的辐亮度轮廓。以法向为太阳中心方向的平面作为DNI参考面，逐方向投影权重和中心余弦为

\[
g_{ik}(\boldsymbol s)=\frac{\boldsymbol n_{ik}\cdot\boldsymbol s}{\boldsymbol s_{0k}\cdot\boldsymbol s},
\qquad c_{ik}=\boldsymbol n_{ik}\cdot\boldsymbol s_{0k}.
\tag{7}
\]
\end{minipage}\par

分母把参考面上的该方向通量还原到传播截面，分子投影到镜面。将光源方向分解为沿太阳中心轴与垂直该轴的分量，式（7）的轴向项为 \(c_{ik}\)，横向项在轴对称方位平均后为零，因而得到 \(\mathbb E[g_{ik}]=c_{ik}\)。

该关系只在完整角域满足镜面正面、地平线上方及参考面投影为正时使用。每个设计均检查全部规定时点的角域适用条件，不截断失败方向后重新归一。式（7）已经包含入射投影，后续不能重复扣除同一余弦。[2,3]

\subsection*{3.4 联合路径事件与三层功率}

从源点向光源回查，若无其他镜面或接收器实体遮阴，则 \(S=1\)。反射路径在规定终点前未被其他镜面截住，则 \(B=1\)。若接收器首次实体接触是有效侧面外入，则 \(R=1\)。有效接收事件为 \(SBR=1\)，重叠障碍只扣能量一次，接收后方的障碍不再削减已经接收的功率。

\par\noindent\begin{minipage}{\linewidth}
未接收的反射光线采用接收器包围球后切平面作为遮挡查询终点。令 \(\boldsymbol y\) 为平面任意点，则

\[
(\boldsymbol y-\boldsymbol c_R)\cdot\boldsymbol t_i=R_b,
\qquad R_b=\sqrt{3.5^2+4^2}.
\tag{8}
\]
\end{minipage}\par

若射线先接触接收器实体，则以首次接触提前终止。所有终止参数必须为正。式（8）是损失分项的记账边界，不是新增实体；题面没有唯一规定未接收射线的终点。改变该终点可能改变阴影遮挡与截断分量而保持净接收不变，因此本文不把两个分项称为唯一的物理分解。

\par\noindent\begin{minipage}{\linewidth}
在镜面面积均匀分布和式（6）的方向分布上取联合期望。令 \(C_{ik}=\rho\,\mathrm{DNI}_kA_i\)，三层功率定义为

\[
\Pi_{0,ik}=C_{ik}c_{ik},\qquad
\Pi_{1,ik}=C_{ik}\mathbb E[g_{ik}SB],\qquad
\Pi_{2,ik}=C_{ik}\mathbb E[g_{ik}SBR].
\tag{9}
\]
\end{minipage}\par

\par\noindent\begin{minipage}{\linewidth}
三层分别为基准反射、联合存活与有效接收功率，均以kW计且尚未乘大气透射。按题给中心距离公式

\[
\tau_i=0.99321-0.0001176d_i+1.97\times10^{-8}d_i^2,
\qquad d_i=\|\boldsymbol c_R-\boldsymbol c_i\|\le1000\,\mathrm m.
\tag{10}
\]
\end{minipage}\par

\par\noindent\begin{minipage}{\linewidth}
在相应分母为正时，五因子关系与单镜功率为

\[
\begin{aligned}
\eta_{{\rm sb},ik}&=\frac{\Pi_{1,ik}}{\Pi_{0,ik}},
&\eta_{{\rm trunc},ik}&=\frac{\Pi_{2,ik}}{\Pi_{1,ik}},\\
\eta_{ik}&=\rho c_{ik}\tau_i\eta_{{\rm sb},ik}\eta_{{\rm trunc},ik},
&P_{ik}&=\tau_i\Pi_{2,ik}=\mathrm{DNI}_kA_i\eta_{ik}.
\end{aligned}\tag{11}
\]
\end{minipage}\par

条件效率使用同一事件体系的分母，不要求阴影与遮挡独立。反射率、投影和大气透射各计一次，不附加接收面余弦、反平方因子或未经提供的热电转换系数。

\subsection*{3.5 汇总对象与异面积关系}

\par\noindent\begin{minipage}{\linewidth}
设月份 \(m\) 的五个时点为 \(\mathcal K_m\)。总功率先逐镜求和，再按时点平均：

\[
\begin{aligned}
A_{\rm tot}&=\sum_iA_i,&P_k&=\mathrm{DNI}_k\sum_iA_i\eta_{ik},\\
\overline P_m&=\frac15\sum_{k\in\mathcal K_m}P_k,
&\overline P&=\frac1{60}\sum_{k=1}^{60}P_k,\\
\overline q_m&=\frac{\overline P_m}{A_{\rm tot}},
&\overline q&=\frac{\overline P}{A_{\rm tot}}.
\end{aligned}\tag{12}
\]
\end{minipage}\par

\par\noindent\begin{minipage}{\linewidth}
对任一已定义效率分量或总效率 \(e_{ik}\)，题表统一采用固定镜数总体的算术平均；面积平均另作辅助统计：

\[
\overline e_{N,m}=\frac1{5N}\sum_{k\in\mathcal K_m}\sum_i e_{ik},
\qquad
\overline e_{A,m}=\frac1{5A_{\rm tot}}\sum_{k\in\mathcal K_m}\sum_iA_i e_{ik}.
\tag{13}
\]
\end{minipage}\par

年度将5改为60并取全部时点。等面积时两种效率平均相同，异面积时一般不同；全场能量总量比 \(\sum\Pi_2/\sum\Pi_1\) 又具有不同分母。总效率先逐镜逐时形成乘积后平均，不以各分量均值乘积代替；式（12）不能改成平均DNI与平均效率的乘积，也不能用未加权的单镜单位面积功率平均代替 \(\overline P/A_{\rm tot}\)。

\section*{4 随机求解与数值确认}

\subsection*{4.1 加权积分及批次池化}

\par\noindent\begin{minipage}{\linewidth}
在镜面面积与光源立体角的乘积域上抽样。四个相互独立的均匀变量 \(U_j\in[0,1)\) 生成

\[
\xi=w_i(U_1-1/2),\quad\zeta=h_i(U_2-1/2),\quad
\mu=1-2\sin^2(\beta/2)U_3,\quad\psi=2\pi U_4.
\tag{14}
\]
\end{minipage}\par

其中 \(\mu=\cos\theta\)。式（14）采用实际实现的随机字映射；以 \(1-U_3\) 替换 \(U_3\) 得到的上升式与之等分布，但不能用于逐射线的相同随机字复放。联合抽样密度为 \(p_\Omega/A_i\)，已经与物理方向份额匹配，贡献中不再重复乘入该密度。

\par\noindent\begin{minipage}{\linewidth}
每镜每时点保存三类加权和、平方和及事件计数。对实际纳入且不重复的源样本，记

\[
M_0=\sum g,\qquad M_1=\sum gSB,\qquad M_2=\sum gSBR.
\tag{15}
\]
\end{minipage}\par

\par\noindent\begin{minipage}{\linewidth}
原始积分 \(M_\ell/n_{\rm tot}\)、自归一比例和条件比是不同估计量。正式点估计在有效状态下使用解析尺度：

\[
\widehat\eta_{\rm sb}=\frac{M_1}{M_0},\quad
\widehat\eta_{\rm trunc}=\frac{M_2}{M_1},\quad
\widehat\eta_{ik}=\rho c_{ik}\tau_i\frac{M_2}{M_0},\quad
\widehat P_{ik}=\mathrm{DNI}_kA_i\widehat\eta_{ik}.
\tag{16}
\]
\end{minipage}\par

先池化各独立批次的原始统计，再计算比例，不平均每批条件比；半量前缀与终值相关，不重复并入样本数。自归一比值在有限样本下可能有偏，解析尺度不能消除全部比值偏差。原始积分另行重建，用于检查估计形式差异和 \(\mathbb E[g]=c\) 的数值对应。事件包含关系给出 \(0\le M_2\le M_1\le M_0\)，不通过裁剪效率取得该关系。

\subsection*{4.2 零值、未知边界与固定总体}

抽样零只表示当前样本没有存活或接收路径，不能自动成为真实零证据。真实零需要当前完整场景的可核验证书；有独立证据确定零接收时，总功率可以合法为零，同时仍保留相应条件效率未定义。输入失败、统计非法和几何边界不确定各自报告，不能伪装成零功率。

对正存活分母下的零接收，保留零接收计数及原始诊断比值，不认定真零；在第2、3问正式汇总中，无证书的该状态仍为数值未决，其截断、总效率和功率保留NA。若固定总体含必需未定义分项，其对应均值也保留NA，不填0或1、不删除镜面或时点。三问最终确认中相关必需状态均已可定义；这不意味任意低样本搜索评分也具备正式含义。

\subsection*{4.3 独立批次与工作指标}

\par\noindent\begin{minipage}{\linewidth}
独立性的统计单位为完整批次。设 \(\mathcal T\) 为从池化原始统计重建某个汇总量的函数，\(\mathcal S_b\) 为第 \(b\) 批覆盖全部对象及时点的统计量，最终确认均采用 \(B_0=8\) 个独立批次。删除一整批后

\[
\widehat Q=\mathcal T\!\left(\sum_b\mathcal S_b\right),\qquad
\widehat Q_{(-b)}=\mathcal T\!\left(\sum_{j\ne b}\mathcal S_j\right),\qquad
\mathrm{SE}_{\rm JK}^2=\frac{B_0-1}{B_0}\sum_b(\widehat Q_{(-b)}-\overline Q_{(-)})^2.
\tag{17}
\]
\end{minipage}\par

式（17）保留批内实际相关结构，不未经说明再除以镜数或时点数的平方根。\(H_Q=t_{7,0.975}\mathrm{SE}_{\rm JK}\) 是近似逐项95\%半宽，依赖独立批次、分母稳定性和刀切近似；不是严格误差界，也不保证多表项同时覆盖。

\par\noindent\begin{minipage}{\linewidth}
三问冻结实验有共同诊断思想，但工作指标不能事后改写为完全相同。第1问以8批独立初始评价I128和8批确认评价C256构成

\[
U_Q^{(1)}=\max\{H_Q,|Q_{\rm C256}-Q_{\rm C128}|,
|Q_{\rm C256}-Q_{\rm I128}|,|Q_{\rm C256}-Q_{\rm raw}|\}.
\tag{18}
\]
\end{minipage}\par

\par\noindent\begin{minipage}{\linewidth}
第2、3问则使用各自冻结确认的终值、相关半量前缀和原始积分，再加入预登记的低存活全范围影响量：

\[
U_Q^{(2,3)}=\max\{H_Q,|\widehat Q-\widehat Q_{1/2}|,
|\widehat Q-\widehat Q_{\rm raw}|\}+D_{{\rm low},Q}.
\tag{19}
\]
\end{minipage}\par

第1问事后补查的低存活影响单列，不回写式（18）。式（18）、（19）都是数值工作指标，不能给整体赋予严格95\%覆盖含义，也不包含物理模型偏差。汇总效率目标为绝对工作指标不超过0.001；总功率和单位面积功率为相对指标不超过0.5\%。对接近零的功率不能直接用相对量认证；最终所需功率汇总未触发该问题。

\par\noindent\begin{minipage}{\linewidth}
冻结设计的额定数值工作判据为

\[
L_P=\widehat{\overline P}-U_{\overline P}\ge60000\,\mathrm{kW},
\tag{20}
\]
\end{minipage}\par

并同时要求几何、模型适用域、必需状态及汇总精度检查通过。它是所采用模型下的工作判据，不能降格成点估计超过额定值，也不能称为工程可靠性的置信保证。

\section*{5 第1问：固定镜场评价}

第1问直接使用附件坐标，统一镜面6m\(\times\)6m、安装高度4m，总面积62820m\textsuperscript{2}。所有1745面镜在60个规定时点形成104700个组合，每个时点完整镜场均参与潜在障碍集合。初始组与确认组每批每组合分别为128和256条源样本，最终点估计只池化确认组的2048条源样本。三问年度结果集中列于表2，月度结果分别在各问给出，以避免重复符号和列名解释。


\par\medskip\noindent\begin{minipage}{\linewidth}
\small
\textbf{表2　三问年样本平均结果（对应题目表2）}\par\smallskip
\begin{tabular}{>{\raggedright\arraybackslash}p{0.145000\dimexpr\linewidth-14\tabcolsep\relax}>{\raggedright\arraybackslash}p{0.120000\dimexpr\linewidth-14\tabcolsep\relax}>{\raggedright\arraybackslash}p{0.120000\dimexpr\linewidth-14\tabcolsep\relax}>{\raggedright\arraybackslash}p{0.120000\dimexpr\linewidth-14\tabcolsep\relax}>{\raggedright\arraybackslash}p{0.130000\dimexpr\linewidth-14\tabcolsep\relax}>{\raggedright\arraybackslash}p{0.135000\dimexpr\linewidth-14\tabcolsep\relax}>{\raggedright\arraybackslash}p{0.230000\dimexpr\linewidth-14\tabcolsep\relax}}
\toprule
问题/设计 & \(\eta\) & \(\eta_{\rm cos}\) & \(\eta_{\rm sb}\) & \(\eta_{\rm trunc}\) & \(\overline P\)（MW） & \(\overline q\)（kW/m\textsuperscript{2}） \\
\midrule
第1问 & 0.5766 & 0.7565 & 0.9279 & 0.9381 & 35.249 & 0.5611 \\
第2问 R027 & 0.4684 & 0.7595 & 0.7733 & 0.9258 & 60.158 & 0.4563 \\
第3问 Q3F013 & 0.4724 & 0.7595 & 0.7789 & 0.9259 & 60.020 & 0.4607 \\
\bottomrule
\end{tabular}
\end{minipage}\par\medskip

表2中后两问分别使用各自正式确认值；第3问的相对改善另见表10的新配对实验。\(\eta,\eta_{\rm cos},\eta_{\rm sb},\eta_{\rm trunc}\) 依次为总光学、余弦、联合阴影遮挡和条件截断的镜数均值，均无量纲。年度总功率换算为MW，\(\overline q\) 单位为kW/m\textsuperscript{2}。


\par\medskip\noindent\begin{minipage}{\linewidth}
\small
\textbf{表3　固定镜场每月21日平均效率及单位面积功率（对应题目表1）}\par\smallskip
\begin{tabular}{>{\raggedright\arraybackslash}p{0.166667\dimexpr\linewidth-12\tabcolsep\relax}>{\raggedright\arraybackslash}p{0.166667\dimexpr\linewidth-12\tabcolsep\relax}>{\raggedright\arraybackslash}p{0.166667\dimexpr\linewidth-12\tabcolsep\relax}>{\raggedright\arraybackslash}p{0.166667\dimexpr\linewidth-12\tabcolsep\relax}>{\raggedright\arraybackslash}p{0.166667\dimexpr\linewidth-12\tabcolsep\relax}>{\raggedright\arraybackslash}p{0.166667\dimexpr\linewidth-12\tabcolsep\relax}}
\toprule
月份 & \(\eta\) & \(\eta_{\rm cos}\) & \(\eta_{\rm sb}\) & \(\eta_{\rm trunc}\) & \(\overline q_m\)（kW/m\textsuperscript{2}） \\
\midrule
1 & 0.5377 & 0.7199 & 0.9133 & 0.9436 & 0.4687 \\
2 & 0.5640 & 0.7404 & 0.9310 & 0.9375 & 0.5318 \\
3 & 0.5830 & 0.7611 & 0.9346 & 0.9350 & 0.5799 \\
4 & 0.6001 & 0.7793 & 0.9360 & 0.9348 & 0.6176 \\
5 & 0.6103 & 0.7893 & 0.9365 & 0.9360 & 0.6378 \\
6 & 0.6137 & 0.7924 & 0.9366 & 0.9366 & 0.6440 \\
7 & 0.6103 & 0.7892 & 0.9365 & 0.9360 & 0.6376 \\
8 & 0.5994 & 0.7786 & 0.9359 & 0.9348 & 0.6162 \\
9 & 0.5822 & 0.7601 & 0.9346 & 0.9352 & 0.5779 \\
10 & 0.5615 & 0.7378 & 0.9299 & 0.9382 & 0.5251 \\
11 & 0.5348 & 0.7182 & 0.9109 & 0.9441 & 0.4623 \\
12 & 0.5219 & 0.7111 & 0.8994 & 0.9457 & 0.4343 \\
\bottomrule
\end{tabular}
\end{minipage}\par\medskip

表3的规定月样本中，单位面积功率在 6 月最高、12 月最低，未舍入值分别为 0.644006 与 0.434324 kW/m\textsuperscript{2}。余弦与总效率的月序变化相近，但该关联不能分离DNI、几何和其他损失的因果贡献。截断分量的极值月份与总效率不同，也不能用单一分量代替全场评价。

第1问年平均光学效率为 0.5766，功率为 35.249 MW；其未达到后续设计题的60MW，不构成本问计算失败。图1展示三项固定结果在规定月样本上的单位面积功率，仅用于比较计算对象与月序变化。第1问的面积和布局不同，较高的单位面积值不说明它满足后两问的额定约束。

\par\medskip\noindent\begin{minipage}{\linewidth}\centering
\includegraphics[width=\linewidth,height=0.72\textheight,keepaspectratio]{\detokenize{图表/图1-三问月度功率.png}}\par
{\small\raggedright 图1　三问规定月样本的单位面积热功率。第2问曲线采用原正式确认值；第3问改善判定不由这两条曲线之差替代。连线只连接12个离散月样本，不表示连续全年积分；未显示误差条，数值工作指标见第8节。\par}\end{minipage}\par\medskip

\section*{6 第2问：统一参数设计}

\subsection*{6.1 目标与几何可行集}

\par\noindent\begin{minipage}{\linewidth}
完整设计为 \(\mathcal D_2=(\boldsymbol a,w,h,z,N,\{\boldsymbol p_i\})\)。面积为 \(A_{\rm tot}=Nwh\)，目标保持为

\[
\max_{\mathcal D_2\in\mathcal F_2}\frac{\overline P(\mathcal D_2)}{Nwh},
\qquad \overline P(\mathcal D_2)\ge60000\,\mathrm{kW}.
\tag{21}
\]
\end{minipage}\par

额定下限是对“达到60MW”的采用解释，不要求各时点均达到额定，也不改成总功率等式。只有总面积固定时最大功率才与原目标等价；只有总功率固定时最小面积才等价，式（21）不满足这两种一般前提。

\par\noindent\begin{minipage}{\linewidth}
场界和禁布区按中心位置判断，固定场地不随塔位平移。几何条件为

\[
\begin{aligned}
&\|\boldsymbol a\|\le350,\quad \|\boldsymbol p_i\|\le350,
\quad\|\boldsymbol p_i-\boldsymbol a\|\ge100,\\
&\|\boldsymbol p_i-\boldsymbol p_j\|\ge w+5\quad(i<j),\\
&2\le w,h\le8,\quad2\le z\le6,\quad z>h/2,\quad N\in\mathbb Z_{>0}.
\end{aligned}\tag{22}
\]
\end{minipage}\par

中心边界是采用的几何解释，不额外要求全部转动镜面始终位于场界内。由于宽边水平，镜面最低点为 \(z-h|v_z|/2\)；全俯仰不触地采用严格正净空 \(z>h/2\)。有限搜索另取 \(w\ge h\)，这是搜索范围限制，不将题面“通常”改成新增硬约束。数值容差只报告近边界对象，不放行实际负间距或零净空。

\subsection*{6.2 圆环与六扇区直边层带表示}

首先以塔为生成中心，用环半径、半开角域、整数计数和相位生成圆环点，之后按固定场界及塔周禁区裁剪。完整圆环的名义计数由弦长不小于目标间距决定；实际跨环间距仍须检查。完整角域和较保守径向间隔限制了该族的密度，不能把其搜索结果推广到任意镜位。

\par\noindent\begin{minipage}{\linewidth}
为扩大表示范围，采用六扇区直边层带。令 \(d_*=w+5+\varepsilon\)，其中 \(\varepsilon\) 为正的数值搜索间隙，公共旋转角为 \(\phi\)。取

\[
\begin{aligned}
\boldsymbol b_s&=d_*(\cos(\phi+s\pi/3),\sin(\phi+s\pi/3)),\quad s=0,\ldots,5,\\
\widetilde{\boldsymbol p}_{\ell sj}
&=\boldsymbol a+(\ell-j)\boldsymbol b_s+j\boldsymbol b_{s+1},
\quad j=0,\ldots,\ell-1,\quad\ell\ge1.
\end{aligned}\tag{23}
\]
\end{minipage}\par

令 \(\boldsymbol b_6=\boldsymbol b_0\)。每条边排除末端、下一边包含该顶点，避免重复槽位；第 \(\ell\) 层有 \(6\ell\) 个互异名义点。同层欧氏半径不全相同，不能把直边层带称为调整圆环半径的同一构造。

\par\noindent\begin{minipage}{\linewidth}
在公共平移、旋转和格距不变时，不同点之差可写成 \(a\boldsymbol b_0+b\boldsymbol b_1\)，其中 \((a,b)\) 为非零整数对，故

\[
\|\Delta\boldsymbol p\|^2=d_*^2(a^2+ab+b^2)\ge d_*^2.
\tag{24}
\]
\end{minipage}\par

式（24）给出包括跨层、跨扇区在内的构造间距依据，但不能直接用于各层独立缩放或旋转。生成足够层数后逐点裁剪，实际镜数按保留清单重新统计；裁剪后另以展开坐标检查全部无序镜对、重复身份、场界及净空，不仅依赖生成器的名义参数。

\subsection*{6.3 受控搜索与可行性修复}

初轮实际比较14个12时点粗筛候选，其中5个接受全60时点较精评价；最好功率候选C22未通过额定工作判据，因此随后在六扇区直边层带附近进行实际修复。原问题目标不变，但在没有可行基线时，先按全60时点功率缩小缺口；出现达到搜索保留条件的候选后，再比较单位面积功率。各次增删镜面或参数改变都展开完整设计，重新建立遮挡关系，不使用镜数单调性或单镜贡献简单加总。

修复实际比较R023—R027五个低样本候选，覆盖宽高、非正方形尺寸和塔位变化；安装高度均为6m，旋转角均为 \(\pi/6\)，数值间隙0.05m。低样本层为2批、每批16条，较精层为4批、每批64条，均覆盖全部镜面及60时点。五个低样本候选存在抽样零，只用于启发式安排；不以其超过60MW的评分认证可行性。


\par\medskip\noindent\begin{minipage}{\linewidth}
\small
\textbf{表4　可行性修复的较精全60时点比较（搜索层）}\par\smallskip
\begin{tabular}{>{\raggedright\arraybackslash}p{0.200000\dimexpr\linewidth-10\tabcolsep\relax}>{\raggedright\arraybackslash}p{0.200000\dimexpr\linewidth-10\tabcolsep\relax}>{\raggedright\arraybackslash}p{0.200000\dimexpr\linewidth-10\tabcolsep\relax}>{\raggedright\arraybackslash}p{0.200000\dimexpr\linewidth-10\tabcolsep\relax}>{\raggedright\arraybackslash}p{0.200000\dimexpr\linewidth-10\tabcolsep\relax}}
\toprule
候选 & 镜数 & 总面积（m\textsuperscript{2}） & 功率评分（kW） & 单位面积评分（kW/m\textsuperscript{2}） \\
\midrule
R023 & 2952 & 134500.5000 & 60506.279 & 0.4498591 \\
R027 & 2981 & 131827.2725 & 60165.183 & 0.4563940 \\
\bottomrule
\end{tabular}
\end{minipage}\par\medskip

表4为较精搜索数据，R023与R027的必需状态均已可定义。两者均通过当时冻结的搜索筛选后，按单位面积目标选取R027，而没有因R023总功率较高就改变目标。R027相对R023的搜索配对差大于该层数值波动诊断，但这不是冻结后的独立确认，也没有校正整个搜索的选择效应。与失败C22相比，最终R027提高了总功率，单位面积功率反而下降，说明可行性修复具有代价。

候选冻结在实际交付精度上，原始坐标按10位小数形成，完整清单按17位有效数字存储。独立确认采用未参与搜索的8批\(\times\)256条；每个时点全部镜面仍是潜在障碍。实际工作簿读回值与冻结清单逐字段一致，确认覆盖2981面镜与60时点的178860个组合。

\subsection*{6.4 最终统一参数与结果}


\par\medskip\noindent\begin{minipage}{\linewidth}
\small
\textbf{表5　R027设计参数（对应题目表3）}\par\smallskip
\begin{tabular}{>{\raggedright\arraybackslash}p{0.200000\dimexpr\linewidth-10\tabcolsep\relax}>{\raggedright\arraybackslash}p{0.200000\dimexpr\linewidth-10\tabcolsep\relax}>{\raggedright\arraybackslash}p{0.200000\dimexpr\linewidth-10\tabcolsep\relax}>{\raggedright\arraybackslash}p{0.200000\dimexpr\linewidth-10\tabcolsep\relax}>{\raggedright\arraybackslash}p{0.200000\dimexpr\linewidth-10\tabcolsep\relax}}
\toprule
塔位（m） & 统一宽\(\times\)高（m） & 统一安装高（m） & 镜数 & 总面积（m\textsuperscript{2}） \\
\midrule
\((0,-15)\) & \(6.65\times6.65\) & 6.00 & 2981 & 131827.2725 \\
\bottomrule
\end{tabular}
\end{minipage}\par\medskip

表5给出的R027总面积为 131827.2725 m\textsuperscript{2}，年度正式值见表2，规定月度结果见表6。最小全对镜间距余量约0.05m，最小全俯仰净空2.675m；这里的几何余量不等于功率的工程裕度。正式逐镜坐标见result2.xlsx。


\par\medskip\noindent\begin{minipage}{\linewidth}
\small
\textbf{表6　R027每月21日平均效率及单位面积功率（对应题目表1）}\par\smallskip
\begin{tabular}{>{\raggedright\arraybackslash}p{0.166667\dimexpr\linewidth-12\tabcolsep\relax}>{\raggedright\arraybackslash}p{0.166667\dimexpr\linewidth-12\tabcolsep\relax}>{\raggedright\arraybackslash}p{0.166667\dimexpr\linewidth-12\tabcolsep\relax}>{\raggedright\arraybackslash}p{0.166667\dimexpr\linewidth-12\tabcolsep\relax}>{\raggedright\arraybackslash}p{0.166667\dimexpr\linewidth-12\tabcolsep\relax}>{\raggedright\arraybackslash}p{0.166667\dimexpr\linewidth-12\tabcolsep\relax}}
\toprule
月份 & \(\eta\) & \(\eta_{\rm cos}\) & \(\eta_{\rm sb}\) & \(\eta_{\rm trunc}\) & \(\overline q_m\)（kW/m\textsuperscript{2}） \\
\midrule
1 & 0.4321 & 0.7284 & 0.7499 & 0.9385 & 0.3773 \\
2 & 0.4558 & 0.7464 & 0.7691 & 0.9301 & 0.4303 \\
3 & 0.4771 & 0.7642 & 0.7843 & 0.9215 & 0.4747 \\
4 & 0.4912 & 0.7792 & 0.7888 & 0.9169 & 0.5055 \\
5 & 0.4985 & 0.7869 & 0.7890 & 0.9163 & 0.5209 \\
6 & 0.5009 & 0.7891 & 0.7892 & 0.9164 & 0.5256 \\
7 & 0.4983 & 0.7868 & 0.7891 & 0.9162 & 0.5206 \\
8 & 0.4906 & 0.7786 & 0.7886 & 0.9169 & 0.5043 \\
9 & 0.4760 & 0.7633 & 0.7835 & 0.9221 & 0.4725 \\
10 & 0.4530 & 0.7441 & 0.7670 & 0.9310 & 0.4242 \\
11 & 0.4298 & 0.7269 & 0.7477 & 0.9395 & 0.3721 \\
12 & 0.4181 & 0.7205 & 0.7338 & 0.9448 & 0.3482 \\
\bottomrule
\end{tabular}
\end{minipage}\par\medskip

冻结设计的年平均功率为 60157.7837 kW，功率工作指标为 4.5146 kW，额定工作下限为 60153.2690 kW，比60MW高 153.2690 kW。其全部必需表项和数值工作目标通过。R027是有限搜索得到的已确认可行设计，不是对任意塔位、自由镜位或连续尺寸空间的全局最优性证明。

\section*{7 第3问：固定布局下的分组异尺寸设计}

\subsection*{7.1 原问题与实际搜索域}

\par\noindent\begin{minipage}{\linewidth}
完整逐镜设计为 \(\mathcal D_3=(\boldsymbol a,N,\{\boldsymbol p_i,w_i,h_i,z_i\})\)，面积和原目标为

\[
A_i=w_ih_i,\quad A_{\rm tot}=\sum_iA_i,\qquad
\max_{\mathcal D_3\in\mathcal F_3}\frac{\overline P(\mathcal D_3)}{A_{\rm tot}},
\quad\overline P(\mathcal D_3)\ge60000\,\mathrm{kW}.
\tag{25}
\]
\end{minipage}\par

\par\noindent\begin{minipage}{\linewidth}
原可行集 \(\mathcal F_3\) 允许改变塔位、镜数、镜位及逐镜参数。逐镜宽高与安装高度范围沿用式（22），全俯仰净空逐镜满足 \(z_i>h_i/2\)。异宽中心间距采用对称解释

\[
\|\boldsymbol p_i-\boldsymbol p_j\|\ge\max(w_i,w_j)+5,\qquad i<j.
\tag{26}
\]
\end{minipage}\par

式（26）同尺寸时还原第2问，是第3问采用的扩展解释，不是三问共同的题面逐字条件。场界及禁区继续按中心位置判断。逐镜 \(w_i\ge h_i\) 仍只作为搜索限制。

实际搜索取 \(\mathcal G_6\subseteq\mathcal F_3\)：固定R027塔位、镜数和全部平面位置，仅在6组共享的宽、高、安装高度中变化。该固定策略不是原问题硬约束。所有镜子取相同参数能表示原R027，也不要求每面镜必须不同。自由度增加只扩大可表示设计集，不保证有限搜索一定改进。

\subsection*{7.2 可复现分组与耦合约束}

\par\noindent\begin{minipage}{\linewidth}
按镜心至塔的平面距离排序，同距按稳定镜身份排序，秩从0开始；分成三个径向秩组，再按相对塔位方位角的两个半开半平面划分。角度归一为 \(\theta_i\in[0,2\pi)\)，则

\[
r_i=\min\!\left(2,\left\lfloor\frac{3\operatorname{rank}_i}{N}\right\rfloor\right),
\quad b_i=\left\lfloor\frac{\theta_i}{\pi}\right\rfloor,
\quad g_i=2r_i+b_i.
\tag{27}
\]
\end{minipage}\par

\par\noindent\begin{minipage}{\linewidth}
组成员在光学评分前固定，尺寸变化不重新分组。固定镜位的宽度上限由全部相邻距离确定：

\[
w_i^{\rm cap}=\min\!\left(8,\min_{j\ne i}\|\boldsymbol p_i-\boldsymbol p_j\|-5\right),
\qquad w_g^{\rm cap}=\min_{i:g_i=g}w_i^{\rm cap}.
\tag{28}
\]
\end{minipage}\par

组宽不得超过该上限。R027的共同上限约6.70m，不能靠显示舍入越过实际边界。改变任一组参数后，逐镜面积、有限边界、三维中心、遮挡筛选及缓存均重新绑定；不能只计算被改镜子的功率，其对其他镜子的阴影与遮挡也属于新场景。

\subsection*{7.3 分层比较与冻结确认}

实际完成14个2批\(\times\)16条的全时点低样本候选，并对4个对象进行4批\(\times\)64条的较精比较，见附录表12。各层均为2981面镜及60时点；相关半量前缀只作诊断，不重复池化。候选冻结前保留R027基线，功率失去可行性或必需状态未解决时不能以较高单位面积评分替代额定条件。

Q3R013在较精对象中单位面积评分最高，扣除搜索工作量后的功率正余量约8.638kW。它低于预先记录的额外30kW软偏好，但仍满足当时冻结的正余量条件；没有临时降低60MW要求，也没有更换种子到通过为止。按实际交付精度冻结为Q3F013后，使用未继续调参的新8批\(\times\)256条进行确认。

候选与R027新基线按稳定位置身份、时点和批次配对，使用共同随机输入；各独立批次之间仍保持独立。对每批删除重建后的 \(\Delta q=\overline q_{\rm cand}-\overline q_{\rm base}\) 计算差值工作指标，而不是先相减功率再除以任一设计面积。旧第2问R027确认数据已用于后续分析，不充当这次新的独立比较数据。

\subsection*{7.4 异尺寸设计及正式结果}


\par\medskip\noindent\begin{minipage}{\linewidth}
\small
\textbf{表7　Q3F013的六组参数}\par\smallskip
\begin{tabular}{>{\raggedright\arraybackslash}p{0.166667\dimexpr\linewidth-12\tabcolsep\relax}>{\raggedright\arraybackslash}p{0.166667\dimexpr\linewidth-12\tabcolsep\relax}>{\raggedright\arraybackslash}p{0.166667\dimexpr\linewidth-12\tabcolsep\relax}>{\raggedright\arraybackslash}p{0.166667\dimexpr\linewidth-12\tabcolsep\relax}>{\raggedright\arraybackslash}p{0.166667\dimexpr\linewidth-12\tabcolsep\relax}>{\raggedright\arraybackslash}p{0.166667\dimexpr\linewidth-12\tabcolsep\relax}}
\toprule
组号 & 镜数 & 宽（m） & 高（m） & 安装高（m） & 组面积（m\textsuperscript{2}） \\
\midrule
0 & 497 & 6.69 & 6.65 & 6.00 & 22110.7845 \\
1 & 497 & 6.69 & 6.65 & 6.00 & 22110.7845 \\
2 & 497 & 6.69 & 6.65 & 6.00 & 22110.7845 \\
3 & 497 & 6.69 & 6.65 & 6.00 & 22110.7845 \\
4 & 582 & 6.69 & 6.05 & 6.00 & 23556.1590 \\
5 & 411 & 6.69 & 6.65 & 6.00 & 18284.7735 \\
\bottomrule
\end{tabular}
\end{minipage}\par\medskip

表7中组4共有582面镜，其镜高为6.05m，其余2399面为6.65m；各镜宽6.69m、安装高度6m。这里的尺寸“高”不同于安装高度，相同安装高度符合第3问允许逐镜取值的条件。图2给出固定镜位上的镜高分布。

\par\medskip\noindent\begin{minipage}{\linewidth}\centering
\includegraphics[width=\linewidth,height=0.72\textheight,keepaspectratio]{\detokenize{图表/图2-异尺寸镜场.png}}\par
{\small\raggedright 图2　Q3F013的镜高分布。颜色表示镜面边长中的高度；全部安装高度均为6m。横纵轴等比例，外圆为固定场界，虚线圆为随塔位移动的中心禁布区，图示不能替代全镜对约束检查。\par}\end{minipage}\par\medskip


\par\medskip\noindent\begin{minipage}{\linewidth}
\small
\textbf{表8　Q3F013设计参数（对应题目表3）}\par\smallskip
\begin{tabular}{>{\raggedright\arraybackslash}p{0.200000\dimexpr\linewidth-10\tabcolsep\relax}>{\raggedright\arraybackslash}p{0.200000\dimexpr\linewidth-10\tabcolsep\relax}>{\raggedright\arraybackslash}p{0.200000\dimexpr\linewidth-10\tabcolsep\relax}>{\raggedright\arraybackslash}p{0.200000\dimexpr\linewidth-10\tabcolsep\relax}>{\raggedright\arraybackslash}p{0.200000\dimexpr\linewidth-10\tabcolsep\relax}}
\toprule
塔位（m） & 统一宽\(\times\)高（m） & 统一安装高（m） & 镜数 & 总面积（m\textsuperscript{2}） \\
\midrule
\((0,-15)\) &  &  & 2981 & 130284.0705 \\
\bottomrule
\end{tabular}
\end{minipage}\par\medskip

表8中统一尺寸及安装高度栏按第3问允许方式留空，不表示逐镜文件缺少参数。result3.xlsx逐行完整保存位置、宽、高、安装高度及稳定身份映射；总面积 130284.0705 m\textsuperscript{2}。年度结果见表2，月度结果见表9。


\par\medskip\noindent\begin{minipage}{\linewidth}
\small
\textbf{表9　Q3F013每月21日平均效率及单位面积功率（对应题目表1）}\par\smallskip
\begin{tabular}{>{\raggedright\arraybackslash}p{0.166667\dimexpr\linewidth-12\tabcolsep\relax}>{\raggedright\arraybackslash}p{0.166667\dimexpr\linewidth-12\tabcolsep\relax}>{\raggedright\arraybackslash}p{0.166667\dimexpr\linewidth-12\tabcolsep\relax}>{\raggedright\arraybackslash}p{0.166667\dimexpr\linewidth-12\tabcolsep\relax}>{\raggedright\arraybackslash}p{0.166667\dimexpr\linewidth-12\tabcolsep\relax}>{\raggedright\arraybackslash}p{0.166667\dimexpr\linewidth-12\tabcolsep\relax}}
\toprule
月份 & \(\eta\) & \(\eta_{\rm cos}\) & \(\eta_{\rm sb}\) & \(\eta_{\rm trunc}\) & \(\overline q_m\)（kW/m\textsuperscript{2}） \\
\midrule
1 & 0.4363 & 0.7284 & 0.7554 & 0.9388 & 0.3810 \\
2 & 0.4600 & 0.7464 & 0.7744 & 0.9306 & 0.4345 \\
3 & 0.4808 & 0.7642 & 0.7895 & 0.9217 & 0.4790 \\
4 & 0.4949 & 0.7792 & 0.7944 & 0.9168 & 0.5102 \\
5 & 0.5021 & 0.7869 & 0.7948 & 0.9161 & 0.5257 \\
6 & 0.5044 & 0.7891 & 0.7951 & 0.9159 & 0.5304 \\
7 & 0.5020 & 0.7868 & 0.7949 & 0.9160 & 0.5255 \\
8 & 0.4944 & 0.7786 & 0.7943 & 0.9171 & 0.5092 \\
9 & 0.4801 & 0.7633 & 0.7892 & 0.9221 & 0.4771 \\
10 & 0.4570 & 0.7441 & 0.7722 & 0.9314 & 0.4282 \\
11 & 0.4340 & 0.7269 & 0.7531 & 0.9398 & 0.3758 \\
12 & 0.4224 & 0.7205 & 0.7396 & 0.9448 & 0.3518 \\
\bottomrule
\end{tabular}
\end{minipage}\par\medskip

第3问年平均功率为 60020.3221 kW，工作指标 7.8759 kW，工作下限 60012.4462 kW。超过额定值的工作余量仅 12.45 kW，约为60MW的 0.0207\%；这通过已设数值规则，却不能支持“工程上具有充足裕度”的判断。最小全对间距余量约0.01m，最小离地净空2.675m，均按实际交付值核对。

\subsection*{7.5 与统一参数设计的新配对比较}


\par\medskip\noindent\begin{minipage}{\linewidth}
\small
\textbf{表10　第3问新配对实验的年度比较}\par\smallskip
\begin{tabular}{>{\raggedright\arraybackslash}p{0.210000\dimexpr\linewidth-10\tabcolsep\relax}>{\raggedright\arraybackslash}p{0.210000\dimexpr\linewidth-10\tabcolsep\relax}>{\raggedright\arraybackslash}p{0.140000\dimexpr\linewidth-10\tabcolsep\relax}>{\raggedright\arraybackslash}p{0.210000\dimexpr\linewidth-10\tabcolsep\relax}>{\raggedright\arraybackslash}p{0.230000\dimexpr\linewidth-10\tabcolsep\relax}}
\toprule
设计 & 功率（kW） & \(U_P\)（kW） & \(\overline q\)（kW/m\textsuperscript{2}） & \(U_q\)（kW/m\textsuperscript{2}） \\
\midrule
R027新基线 & 60159.7717 & 8.1620 & 0.4563530 & \(6.19141\times10^{-5}\) \\
Q3F013 & 60020.3221 & 7.8759 & 0.4606881 & \(6.04518\times10^{-5}\) \\
差（候选\(-\)基线） & -139.4496 & — & 0.0043351 & \(2.63993\times10^{-5}\) \\
\bottomrule
\end{tabular}
\end{minipage}\par\medskip

表10的R027来自第3问的新配对实验，与表2第2问原确认是同一物理设计、不同随机样本。其功率点估计相差 1.9880 kW，不能为统一数字而覆盖原确认，也不能混用两次实验计算提升。

同一配对数据给出单位面积差 0.0043351 kW/m\textsuperscript{2}，差值工作指标 \(2.63993\times10^{-5}\) kW/m\textsuperscript{2}，扣减后仍为 0.0043087 kW/m\textsuperscript{2}，因此支持所用模型下约 0.94995\% 的单位面积改善。与此同时，总面积减少 1543.202 m\textsuperscript{2}，总功率下降 139.4496 kW。该取舍支持面积利用效率改善，不能改写成总功率同步提高。

第3问年平均面积加权总效率为 0.472923，镜数平均为 0.472360，说明两者在此异面积设计中确有区别。正式题表保持镜数均值，功率按式（12）的真实面积求和，未为取得较高表值切换权重。

\section*{8 验证、数值证据与适用边界}

\subsection*{8.1 四类证据及其作用}

人工解析构型检查有限矩形、圆柱首次接触、反射方向、重叠障碍及接收后方障碍。第1问局部半遮阴构型的解析份额为0.5，随机估计0.501403，确定性求积0.5；有限接收构型解析份额约0.424264，随机估计0.423073。确定性求积曾随节点细化非单调变化，后续解析单元包络包含真值。这些证据支持所测构型，不是任意镜场几何的严格误差证明。

筛选测试对预先按几何条件选定的困难组合，使用相同射线与全对象穷举比较阴影、反射遮挡、接收和未知状态。第1问12个压力组合一致，后续参数化、实际规模及异尺寸实现另作限定回归；第3问记录的是18项限定原型检查，其中包含旧场景与真实异尺寸组合，不能全部称为纯解析人工试验。

独立随机批次提供式（18）、（19）的数值诊断；从保存统计重建题表以及清单读回检查提供数据可复现证据。三问定义有效处五因子与直接能量路径的最大差约 \(3.33\times10^{-16}\)，面积—功率关系差约 \(1.46\times10^{-11}\,\mathrm{kW}\)，属于数据和代数链路核对，不能独立验证真实物理模型。尚未量化的物理偏差是第四类问题，不能由前述文件一致或随机稳定性替代。


\par\medskip\noindent\begin{minipage}{\linewidth}
\small
\textbf{表11　已保存正式确认的汇总数值检查}\par\smallskip
\begin{tabular}{>{\raggedright\arraybackslash}p{0.200000\dimexpr\linewidth-10\tabcolsep\relax}>{\raggedright\arraybackslash}p{0.200000\dimexpr\linewidth-10\tabcolsep\relax}>{\raggedright\arraybackslash}p{0.200000\dimexpr\linewidth-10\tabcolsep\relax}>{\raggedright\arraybackslash}p{0.200000\dimexpr\linewidth-10\tabcolsep\relax}>{\raggedright\arraybackslash}p{0.200000\dimexpr\linewidth-10\tabcolsep\relax}}
\toprule
问题 & 组合数 & 最大效率指标 & 最大功率相对指标（\%） & 低存活组合数 \\
\midrule
Q1 & 104700 & 0.0002182 & 0.03184 & 5 \\
Q2 & 178860 & 0.0002768 & 0.05474 & 2 \\
Q3 & 178860 & 0.0003575 & 0.04743 & 2 \\
\bottomrule
\end{tabular}
\end{minipage}\par\medskip

表11中的工作指标均针对当前模型的数值评价，不是物理误差。第1问低存活补查与冻结指标分别列示；第2、3问低存活影响已经进入各自工作指标，不能重复添加。三问最终必需状态可定义，第2、3问均无池化抽样零、边界未知或非法统计；这一结论不扩展到所有搜索层记录。

\subsection*{8.2 低存活样本的局部限制}

合并存活数少于100只是筛查门槛，不是单镜精度判据，也不等同于加权有效样本量。第1问104700个组合的最小存活数为11，筛出5个组合；第2问及第3问最终确认各有2个。第3问两组合的存活/接收数为39/36和62/62。即使样本接收比等于1，也不能证明真实截断效率为1。

\par\noindent\begin{minipage}{\linewidth}
以 \(f_1,f_2\) 表示联合存活与有效接收份额，保持 \(0\le f_2\le f_1\le1\)，正分母下 \(\eta_{\rm trunc}=f_2/f_1\)。对筛出集合 \(\mathcal L\)，每月或年度时点数为 \(K\)，逐坐标可行范围直径给出保守影响量

\[
\begin{aligned}
D_{{\rm low},\eta_{\rm sb}},D_{{\rm low},\eta_{\rm trunc}}&\le\frac{|\mathcal L|}{KN},\\
D_{{\rm low},\eta}&\le\frac1{KN}\sum_{(i,k)\in\mathcal L}\rho c_{ik}\tau_i,\\
D_{{\rm low},P}&\le\frac1K\sum_{(i,k)\in\mathcal L}\mathrm{DNI}_kA_i\rho c_{ik}\tau_i,
\quad D_{{\rm low},q}=D_{{\rm low},P}/A_{\rm tot}.
\end{aligned}\tag{29}
\]
\end{minipage}\par

面积平均效率另用 \(A_i/A_{\rm tot}\) 权重。式（29）不把相互矛盾的分量极值当成同一条光线路径，其逐坐标外包界可能偏宽；零分母处仍未定义。第1问事后逐项叠加诊断的最大效率为 0.0004095，功率相对最大量为 0.0695\%，仍低于工作目标。这不是严格联合置信界，也不证明局部条件比已精确。

\subsection*{8.3 分项约定与物理、搜索局限}

有限终点约定的人工测试表明，改变未接收射线终点可使阴影遮挡效率从1变为0.5、截断效率从0.5变为1，而净接收仍为0.5。这说明分项归属依赖路径边界，不能因乘积一致就把两个分项视为唯一物理定义。三问统一保留这一记账约定，没有根据预期结果调整终点。

角分布、代表日历、DNI公式、反射率示例、理想镜面、有效接收面和省略实体均限定了结果的物理意义。现有资料不能量化所在地实际光源、坡度与跟踪误差、支架和塔身遮挡及接收器热损失。增加随机射线只能减少给定模型的数值波动，不能消除这些偏差；第3问很小的额定工作余量尤其不宜延伸为工程裕度。

第2问只探索有界圆环和直边层带族，第3问固定R027的塔位、镜数与镜位，仅释放6组参数，遗漏自由位置、组内单镜差异及布局密度的联合调整。共同随机数与冻结后确认减少了比较被随机高估的风险，却不证明搜索充分、校正所有选择效应或建立全局最优性。

\section*{9 全文结论}

有限镜面、方向份额投影和联合路径事件构成了三问一致的评价接口。固定镜场得到年样本平均热功率 35.249 MW；统一参数设计R027通过额定数值工作判据，年平均功率 60.158 MW、单位面积功率 0.4563379 kW/m\textsuperscript{2}。分组异尺寸设计Q3F013在固定镜位上减少面积，并在同一新配对实验中以总功率下降为代价，使单位面积功率提高约 0.94995\%。

两个设计的正式清单与冻结确认输入一致，必需表项均满足所采用的汇总数值工作目标。Q3F013的额定工作余量仅 12.45 kW，应与其有限搜索范围及未量化物理偏差一同理解。本文得到的是当前模型下可复现、经独立数值确认的有限搜索设计，不将其解释为连续全年实际能量、充分工程裕度或原问题全局最优。

\section*{资料说明}

[1] 题目及配套附件，本地提供版本。题目物理页1—4；太阳位置、辐照度及光学效率关系见附录，镜位与结果模板见附件。

[2] 本地整理方法卡M01—M05，固定版本v001：输入假设、空间方向反射与有限几何、条件能量比与分布归一化、平均权重与非线性汇总、数值误差验证与复现。配套Datawhale材料限于D01—D03选段，本文引用本地整理与选段的实际阅读范围。

[3] 本地补充资料P01《角尺度来源与边界》和P02《方向分布与待选假设》，固定版本v001。角尺度来源形式为本地转述的工程示例，未直接核对上游原件。布局关系及本文问题专用推导依据题面、基本数学与自身实现。

\section*{附录A 数值实验身份与搜索补充}

第1问正式确认采用根种子2026090402、命名空间912；第2问R027原确认采用2026090505、1990；第3问Q3F013与新R027配对采用2026090506、3190。均为PCG64、8批、每批256条，128条为相关前缀。随机流还绑定规定时点、稳定对象身份及批号；第2问原确认与第3问新基线互不合并。完整命令、版本哈希、原始分块和失败恢复另存可复现说明。


\par\medskip\noindent\begin{minipage}{\linewidth}
\small
\textbf{表12　第3问较精全60时点比较（确认前）}\par\smallskip
\begin{tabular}{>{\raggedright\arraybackslash}p{0.200000\dimexpr\linewidth-10\tabcolsep\relax}>{\raggedright\arraybackslash}p{0.200000\dimexpr\linewidth-10\tabcolsep\relax}>{\raggedright\arraybackslash}p{0.200000\dimexpr\linewidth-10\tabcolsep\relax}>{\raggedright\arraybackslash}p{0.200000\dimexpr\linewidth-10\tabcolsep\relax}>{\raggedright\arraybackslash}p{0.200000\dimexpr\linewidth-10\tabcolsep\relax}}
\toprule
候选 & 总面积（m\textsuperscript{2}） & 功率评分（kW） & 单位面积评分（kW/m\textsuperscript{2}） & 功率搜索余量（kW） \\
\midrule
Q3R000 & 131827.2725 & 60162.8604 & 0.4563764 & 148.1191 \\
Q3R010 & 130737.8667 & 60058.8476 & 0.4593837 & 42.2591 \\
Q3R012 & 130673.4285 & 60079.7114 & 0.4597699 & 63.0714 \\
Q3R013 & 130284.0705 & 60024.4440 & 0.4607197 & 8.6382 \\
\bottomrule
\end{tabular}
\end{minipage}\par\medskip

表12为确认前搜索层结果，Q3R010、Q3R012未接受最终独立确认，不能由本表将其标为已确认可行设计。R027原正式结果、搜索层R027和新配对R027是同一设计在不同实验中的记录，各有用途。较长低样本轨迹、局部异常、源样本身份与代码接口不重复写入正文，均保留于证据文件。

\section*{附录B 结果与可复现文件}

月度表3、6、9分别对应三问的题目表1；表2按行保留三问题目表2的字段与单位；表5、8分别对应后两问的题目表3。实际逐镜清单为result2.xlsx和result3.xlsx。第3问表3统一参数栏为空，但逐镜文件、表7与总面积相互对应。

全文结果由四类有效确认来源生成：第1问原确认、第2问原确认、第3问候选确认及其新配对基线。来源表另外绑定配对差文件，防止使用不同实验拼接改善比例。保存统计的独立重建、工作簿读回和论文表格一致性是可复现检查，不构成新的物理实验。数据、代码、图表入口、局限及实验边界审计随完整论文提供，较大的原始统计以工作区相对路径和校验信息定位。
\end{document}
